% Transcription — fonds Alexandre Grothendieck, shelfmark 64, batch 2 (pages 21-40).
% Established from the facsimile archives/batches/64/batch-02.pdf (a copy of 64.pdf
% fetched through the project relay, no cover sheet: PDF page k = archivists' page 20(N-1)+k),
% per the skill transcribe-grothendieck. The folder is 156 pages in eight batches.
%
% Pass: Opus 5.5 (claude-opus-5-5), 2026-09-27 — first pass, unchecked against the pages by a human.
%
% Pages skipped: 40 (an unrelated typed administrative letter, no mathematics).
% His pagination: he numbers the even archive pages only in this run: p. 22 = his 11,
%   24 = 12, 26 = 13, 28 = 14, 30 = 15, 32 = 16, 34 = 17, 36 = 18. Page 21 continues
%   his §3 (subgroups of S_4, items g)-i)), whose start lies in batch 1's undrafted pp. 14-20.
% What the batch is about: end of the census of subgroups of S_I (p. 21); §4 the combinatorial
%   cube defined by Circ(I) -> J, u -> u^2 (pp. 22-24); §5 free rank-2 modules M over
%   Z/4Z: M_0 = M/2M, the End(M_0)-torsor P_M, End(M_0) = M_0 + F_2^omega, the theorem
%   M <-> (M_0, I_M, S_M) <-> (I, Q, k), GL_Lambda(M) as a fibre product of S_I and the
%   cube group D_Q over {+-1}, the subgroup lattice (28) (pp. 25-31); §6 the unoriented
%   octahedron Delta_M = P(M)(Z/4Z) over J = P(M)(F_2), Aut_oct(Delta) = GL(M)/{+-1},
%   the three two-element sets C_Delta, omega_Delta, Or(Delta) (pp. 32-37); a partial
%   recapitulation (p. 38); a landscape sheet in paler ink over a base S (p. 39).
% Notation fixed once for the batch: \mathfrak{S}_I, \mathfrak{D}_Q, \underline{\omega},
%   \underline{a}, \mathfrak{z} = {1, -id_M} (his cursive z), \Lambda = Z/4Z, M^{\#} for the
%   elements of M not in 2M (his star/hash sign), \mathfrak{gl}/\mathfrak{sl}(M_0) for two
%   cursive symbols he distinguishes, \operatorname{Or}(\Delta), \mathfrak{P}_{2,2}.
% \ill{}: 95. \uncertain{}: 25.

% A shared scratchpad file was overwritten mid-pass by another parallel pass;
%   pages 21-31 were restored from the pass's own record of its readings, not
%   re-read from the tiles a second time. Worth a closer check against the pages.

\documentclass[11pt,a4paper]{article}
\input{../preamble/grothendieck}

\folder{64}
\batch{2}
\pages{21}{40}
\dating{[à partir de 1980- 1982]}
\watermark{Édition de démonstration}
\foldertitle{Modules courbes elliptiques en caractéristiques ≠2 : notes manuscrites (s.d.)}

\begin{document}



\page{21}
\textbf{B)} Considérons ensuite les sous-groupes \add{de $G$} d'ordre multiple de $3$, donc d'ordre $3$, $6$ ou $12$.

\begin{itemize}
  \item[g)] Sous-groupes d'ordre $3$, ce sont les $4$ $3$-sous-groupes de Sylow, indexés par les éléments de $I$
    \marginal{$I$ \quad $4$}
  \item[h)] Les sous-groupes d'ordre $6$ ne sont pas cycliques (il n'y a pas d'éléments d'ordre $6$) donc sont isomorphes : $\mathbb{D}_3$, \struck{ou} ce sont les quatre stabilisateurs des pts de $I$, i.e. les normalisateurs des $3$-sous-groupes de Sylow
    \marginal{$I$ \quad $4$}
  \item[i)] Il y a un seul sous-groupe d'ordre $12$ i.e. d'indice $2$, du fait que \struck{$\mathfrak{S}$} $G_{\mathrm{ab}} \simeq$ \uncertain{$\mathbb{F}_2$}, et c'est $G^+ \simeq \mathfrak{S}_I^+$, ss-groupe \emph{alterné} de $\mathfrak{S}_I$.
    \marginal{$1$}
\end{itemize}
\note{les chiffres de la marge gauche (« $I$ 4 », « $I$ 4 », « 1 ») sont le décompte que l'auteur tient en regard de chaque cas : l'ensemble qui indexe la classe de conjugaison et son cardinal. Dans i), la formule après « du fait que » est surchargée : un $\mathfrak{S}$ biffé, puis $G_{\mathrm{ab}}$, et un dernier symbole barré qui se lit mal. La page~21 poursuit un inventaire commencé dans des pages que le premier lot n'a pas encore transcrites (pp.~14--20).}

Il y a donc \add{exactement} $17$ sous-groupes qui sont des deux-sous-groupes, $\neq e$, se répartissant en six classes de conjugaison (correspondant \uncertain{bien} à $6$ classes d'isomorphismes d'espaces homogènes sous $G$), et $9$ sous-groupes \add{$\neq G$} qui sont d'ordre multiple de $3$, se répartissant en trois autres classes de conjugaison. Cela fait en tout $28$ sous-groupes, donc $30$ en comptant $\{e\}$ et $G$, se répartissant en $9$ classes de conjugaison (\uncertain{resp.} $11$ en comptant $\{e\}$ et $G$), qui correspondent aux classes d'isomorphie d'espaces homogènes sous $G$.
\note{le « 17 » est net ; or $17 + 9 = 26$, non $28$. Le $\mathfrak{S}_4$ a $19$ sous-groupes d'ordre une puissance de $2$ distincts de $\{e\}$, et c'est $19 + 9 = 28$ qui donne le total écrit ensuite. Transcrit tel quel.}

\begin{quote}
Pour bien faire, il faudrait les \uncertain{relations} \ill{} dans un diagramme d'inclusions, ou diagr. de morphismes d'espaces homogènes.
\end{quote}
\note{ces deux dernières lignes, d'une écriture plus petite, sont serrées au bas de la page, à droite.}

\section{4. Relations avec le cube orienté}

\page{22}\note{p.~11 de l'auteur, chiffrée en tête de page. La section~3, dont la page~21 donne la fin, commence dans les pages du premier lot qui ne sont pas encore transcrites.}
Si $I$ est un ens. à $4$ éléments, $J \simeq \mathfrak{P}_{2,2}(I) \simeq V_I^*$ l'ens. à trois éléments associé, on considère l'application de degré $2$
\[
  (*) \qquad
  \underset{\substack{\text{ens. des}\\ \text{permutations}\\ \text{circulaires}\\ \text{sur } I}}{\mathrm{Circ}(I)} \longrightarrow J, \qquad u \longmapsto u^2 .
\]
\note{le $J$ d'arrivée est écrit sur un autre symbole.}

[\emph{NB} \struck{\ill{}} On a aussi l'application de degré $2$
\[
  \underset{\substack{\simeq\ \text{ens. des}\\ \text{« transpositions »}\\ \text{de } I\\ \text{(ou « réflexions »)}}}{\mathfrak{P}_2(I)} \longrightarrow J,
  \qquad A \longmapsto \{A, I \setminus A\}
\]
\[
  \text{ou}\quad \forall\, \tau \text{ réflexion}, \ \exists!\, \tau' \text{ réflexion telle que } \tau\tau' \in V^*
\]
ce sont maintenant des revêtements de \uncertain{nature} distincte, les espaces homogènes \add{sous $G$} $\mathrm{Circ}(I)$ et $\mathfrak{P}_2(I)$ n'étant pas isomorphes (le stabilisateur d'un élément de $\mathrm{Circ}$ est $\simeq \mathbb{Z}/4\mathbb{Z}$, celui d'un élément de $\mathfrak{P}_2(I)$ est $\simeq \mathbb{F}_2 \times \mathbb{F}_2$) --- c'est le premier \add{\struck{$\mathfrak{P}$} $\mathrm{Circ}(I)$}, qui peut s'interpréter comme \struck{l'application} \add{l'ensemble} \ill{} \ill{} d'un cube\ldots].
\note{le « pour » écrit sous « $\tau\tau' \in V^*$ » renvoie, par un trait, à la ligne suivante. Devant « d'un cube », deux mots (« des \ldots\ ») ne sont pas lus.}

L'application de degré $2$ $(*)$ définit \add{\uncertain{bien sûr}} un cube \add{comb.}, dont \uncertain{l'}ens. des faces est $\mathrm{Circ}(I)$, \struck{sommets} \add{(\uncertain{compl.} des sommets)} l'un des \struck{arêtes} (\ill{} l'un des $12$ \ldots) \struck{\ill{}} sections partielles \struck{\ill{}} sur des parties de card.~$2$ de $J$ (resp. en \struck{\ill{}} $8$ sections de $\mathrm{Circ}(I)$ sur $J$), l'incidence entre faces étant le prolongement des \struck{\ill{}} \add{sections} partielles. \struck{À \ill{}} Il y aurait
\marginal{Bien sûr, $G$ opère sur le cube par automorphismes, \ill{} avec l'\ill{} $u \mapsto u^{-1}$}
\note{ce paragraphe est surchargé d'ajouts interlinéaires et de biffures, et son ordre de lecture n'est pas sûr. Le sens mathématique est clair : le cube combinatoire défini par $(*)$ a pour faces les six éléments de $\mathrm{Circ}(I)$, pour sommets les huit sections de $\mathrm{Circ}(I) \to J$, pour arêtes les douze sections partielles au-dessus des parties à deux éléments de $J$, l'incidence étant le prolongement des sections partielles. La note marginale, écrite de biais dans la marge gauche et reliée par une accolade à « un cube comb. », n'est lue qu'en partie.}

\page{23}
lieu d'expliciter
\begin{itemize}
  \item[a)] La correspondance entre \add{les $8$} sommets du cube, et les \add{$8$} éléments d'ordre $3$ de $G$ (« rotations-sommets »)
  \item[b)] La \struck{\ill{}} \add{corr.} entre les $6$ paires d'arêtes antipodiques du cube, et les $6$ \struck{réfl} réflexions dans $G \simeq \mathfrak{S}_I$
  \item[c)] La corr. entre les $6$ \struck{orientations de $\mathrm{Circ}(I)$, et} \add{faces du cube, et} les \struck{\ill{} éléments de $G$} \add{permutations} circulaires (i.e. d'ordre $4$ de $G \simeq \operatorname{Aut}^+(\text{Cube})$) correspondants,
\end{itemize}
et de décrire également une \emph{orientation} \emph{canonique} du cube, ou ce qui revient au même, \ill{} \struck{une trivialisation} iso.
\[
  \underset{\underset{\textstyle \underline{\omega}_J\,(\simeq \underline{\omega}_I)}{\| \mathrm{def}}}{\underline{\omega}_Q}
  \ \simeq\
  \bigwedge_{i \in J}
  \underset{\substack{\text{ens. des}\\ \text{deux rotations, d'ordre 4}\\ \text{de la structure carrée}\\ Q_i \text{ définie par } i \in J\\ \text{i.e. ens. des deux orientations}\\ \text{de } Q_i}}{\underline{\omega}_i}
\]
\note{l'indice de $\underline{\omega}$ à gauche est surchargé ; la lecture $Q$ est celle du lot précédent ($\underline{\omega}_Q$). Sous $\underline{\omega}_i$, « rotations » est ajouté au-dessus d'un mot biffé illisible, et n'est pas sûr ; de même « carrée ».}

Comme on a
\[
  \underline{\omega}_i \simeq
  \underset{\substack{\text{ens.}\\ \text{des deux}\\ \text{codiagonales}\\ \text{de } Q_i}}{\underline{\omega}} \wedge
  \underset{\substack{\text{ens. des diag.}\\ \text{pour le carré } Q_i}}{D_i}, \quad \text{i.e. } D_i = I/\underset{\substack{\text{translation}\\ \text{indexée par } i}}{a_i}
\]
la formule \struck{\ill{}} s'écrit aussi
\[
  \underset{\text{torseur sous } \mathbb{F}_2^J}{\bigwedge_{i \in J} D_i} \simeq \mathbb{1}
\]
Or on a \struck{la suite exacte} l'inclusion
\[
  \underset{\text{torseur sous } V_J(\mathbb{F}_2)}{I} \hookrightarrow \underset{\text{torseur sous } \mathbb{F}_2^J}{\prod_{i \in J} D_i}
\]
\note{les légendes sous $\underline{\omega}$ et $D_i$ sont écrites en petit, en partie l'une sur l'autre ; celle de gauche (« ens.\ des deux codiagonales de $Q_i$ ») est d'une lecture incertaine. Le symbole après $I/$ est surchargé ; il est lu $a_i$ d'après la légende « translation indexée par $i$ ». Sous $\bigwedge D_i$, la légende commence par un mot biffé. Sous les deux termes de l'inclusion, une flèche d'inclusion relie aussi les deux légendes.}

\page{24}\note{p.~12 de l'auteur.}
qui fait une restriction du groupe structural à $V_J(\mathbb{F}_2)$, d'où $\bigwedge_{i \in J} D_i \simeq \mathbb{1}$, OK.

Il y aurait lieu aussi d'interpréter les $10$ classes d'isomorphie d'espaces homogènes $\to(1)$ sous $G \simeq \operatorname{Aut}(\text{Cube})$ (correspondant aux \struck{classes} \add{sous}-groupes d'isotropie $\neq G$, définis à conjugaison près) en termes d'\struck{objets} éléments de structure invariants du cube (tels sommets, arêtes, faces, ou \struck{tels} objets tant antipodiques --- i.e. que le cube garde --- ou \add{$n$-paires} \struck{\ill{}} \add{à \uncertain{iso} près}, déjà explicités), de cardinaux
\[
  8,\ 12,\ 6,\ 4,\ 6,\ 3,\ 24 .
\]
Il y en a un de cardinal $2$, quotient de l'un des sommets (correspondant à l'ens. $\underline{\omega}_I$), un \ill{} d'ordre \add{(ou codiagonales, \ill{})} correspondant aux $3$ diagonales de faces (les \struck{\ill{} d'ordre} \add{\ill{}} \ill{} d'ordre $2$ \ill{} dans $V$), enfin un \struck{\ill{}} $G$ qui \struck{\ill{}} \ill{} des paires d'éléments distincts de $J$ (i.e. \ill{} de $J$), $J$ étant interprété comme l'ens. des « directions de \struck{\ill{}} \add{faces} » \struck{(\ill{})} (ens. des faces$/$antipodisme) --- ou directions d'arêtes. Le compte y est !
\marginal{\uncertain{partition} \ill{} relation \ill{} sommets : dist.~$1$, \ill{} \uncertain{encore} \ill{} arêtes : dist.~$2$ \ldots\ dist.~$3$}
\note{le renvoi « $\to(1)$ » n'a pas de note correspondante sur la page. La liste des cardinaux correspond vraisemblablement aux sommets, arêtes, faces, diagonales, paires d'arêtes antipodiques, paires de faces opposées et repères. Les lignes qui suivent sont d'une lecture très incertaine ; le sens est d'achever l'inventaire des espaces homogènes (cardinaux $2$, $3$, \ldots) et de constater que le compte des dix classes est juste. La note marginale, écrite obliquement dans la marge gauche et en partie barrée d'un trait, n'est lue que par fragments. Au bas de la page, un signe isolé, non lu.}

\section{5. Modules libres de rang 2 sur $\Lambda = \mathbb{Z}/4\mathbb{Z}$.}

\page{25}
Soit $M$ un tel module, et soit
\[
  (1) \qquad M_0 = M/2M = M \otimes_\Lambda \mathbb{F}_2
\]
qui est un vectoriel de rang $2$ sur $\mathbb{F}_2$. Soit $J = M_0^*$, on utilisera la notation indicielle pour les éléments de $M_0^*$, notés $x_i$ ($i \in J$). \struck{Pour tout $i \in J$, soit}
On a une suite exacte de \struck{groupes commutatifs} \add{$\Lambda$-modules}
(2)


\begin{tikzcd}
    0 \arrow[r] & M_0 \arrow[r] & M \arrow[r, "\mathrm{can}"] & M_0 \arrow[r] & 0 \\
    & & M \arrow[ul, dashed, "\mathrm{can}"] \arrow[u, dashed, "2\cdot"'] & &
  \end{tikzcd}

donc $M$ apparaît comme extension de $M_0$ par $M_0$ (extension par $\Lambda$-modules). Pour tout $x \in M_0$, soit $M_x$ son image inverse dans $M$, qui est un torseur sous $M_0$ (i.e. un ens. à quatre éléments, \add{$M_x$}, muni d'un iso. $M_0 \simeq V_{M_x}$).
\note{dans le diagramme (2), les deux flèches pointillées partent d'un $M$ écrit sous la suite : « can » monte vers le premier $M_0$, « $2\cdot$ » vers le $M$ du milieu ; le diagramme est redessiné.}

On posera, pour $i \in J$
\[
  (3) \qquad
  \left\{
  \begin{aligned}
    M_i &= M_{x_i} = \ldots \\
    M_J &= \prod_{i \in J} M_i, \quad \text{un torseur sous } M_0^J .
  \end{aligned}
  \right.
\]
\note{après $M_i = M_{x_i} =$, l'auteur a biffé « image inverse de $\{x_i\}$ » et écrit au-dessus « fibre de $M$ au-dessus de $x_i \in M_0$, $M_0$-torseur » ; devant $M_0^J$, un $\mathbb{F}$ biffé.}

Comme on a
\[
  (4) \qquad \sum_{i \in J} x_i = 0
\]
la structure de groupe sur $M$ donne un iso. can.
\[
  (6) \qquad \bigwedge_{i \in J} M_i = 0
\]
en d'autres termes, introduisant
\[
  (7) \qquad V_J(M_0) = \Bigl\{ (\xi_i)_{i \in J} \in M_0^J \Bigm| \sum \xi_i = 0 \Bigr\}
\]
s'insérant dans la suite exacte
\[
  (8) \qquad 0 \to V_J(M_0) \to M_0^J \to M_0 \to 0 ,
\]
le torseur $M_J$ sous $M_0^J$ provient canoniquement d'un torseur sous $V_J(M_0)$, que je noterai $P_M$ :
\note{il n'y a pas de formule (5) ; le « (6) » est net.}

\page{26}\note{p.~13 de l'auteur, chiffrée « (13) » en tête de page.}
De façon précise, \add{la donnée} des \add{$M_0$-torseurs} $M_i$ et \add{de} l'isom. can. $(6)$ revient à la donnée du torseur $P_M$ sous $V_J(M_0)$, avec
\[
  (9) \qquad M_i \simeq P_M \wedge^{M_0^J} (M_0, \mathrm{pr}_i) \qquad \text{pour } i \in J
\]
où
\[
  \mathrm{pr}_i : M_0^J \to M_0 \qquad \text{la projection d'indice } i \in J .
\]
[Notons qu'on a des iso. canoniques (valables pour \uncertain{tout} vectoriel de rang $2$ sur $\mathbb{F}_2$)]
\[
  (10) \qquad \operatorname{End}(M_0) \simeq V_J(M_0) \xrightarrow{\;\sim\;} M_0 \otimes M_0 \simeq M_0 \oplus \mathbb{F}_2^{\underline{\omega}_{M_0}}
\]
\[
  \underline{\omega}_{M_0} \simeq \underline{\omega}(\underset{J}{\underbrace{M_0^*}}) \quad \text{ens. des orientations de } J .
\]
\[
  (12) \qquad
  \left\{
  \begin{aligned}
    \operatorname{End}(M_0) &\xrightarrow{\;\sim\;} V_J(M_0) \\
    u &\longmapsto (u(x_i))_{i \in J}
  \end{aligned}
  \right.
\]
\[
  \operatorname{End}(M_0) \simeq M_0^\vee \otimes M_0 \simeq M_0 \otimes M_0 \quad \text{car } M_0^\vee \simeq M_0
\]
(valable pour \uncertain{tout} vect. de rang $2$ sur $\mathbb{F}_2$)
\marginal{\emph{NB} Pour \uncertain{tout} vectoriel $W$ sur $\mathbb{F}_2$, on a $(11)$ $\operatorname{Hom}_{\mathbb{F}_2}(M_0, W) \simeq V_J(W)$, car $M_0 \simeq \mathbb{F}_2^J/\mathrm{diag}$, d'où $(12)$.}
\note{la note marginale est écrite obliquement dans la marge gauche, en face de (10)--(12), et une flèche la relie à (12). Le numéro « (10) » est surchargé à l'encre plus noire. La lettre $W$ est lue pour un symbole qui ressemble à un $\omega$ ; le quotient « $/\mathrm{diag}$ » n'est pas sûr.}

\[
  (13) \qquad
  \begin{array}{ccc}
    M_0 & \hookrightarrow & \operatorname{End}(M_0) \\
    \cup & & \cup \\
    M_0^* & \longleftrightarrow & \operatorname{Aut}(M_0) \simeq \mathfrak{S}_J \\
    x_i & \longmapsto & \sigma_i
  \end{array}
\]
où $x_i \mapsto$ \struck{transposition} \add{$\sigma_i$} réflexion de $M_0$ qui fixe $x_i$ (ou \struck{transposition} \add{\ill{}} de $x_j$ et $x_k$)
[\emph{NB} Comme $\sum_{i \in J} \sigma_i = 0$, cela correspond à un hom. $M_0 \to \operatorname{End}(M_0)$]
\note{le $M_0^*$ de la deuxième ligne de (13) est écrit sur un autre symbole ; les signes « $\cup$ » sont ceux de l'auteur (inclusions dressées verticalement).}

\[
  (14) \qquad
  \begin{array}{ccc}
    \mathbb{F}_2^{\underline{\omega}_{M_0}} & \hookrightarrow & \operatorname{End}(M_0) \\
    \cup & & \cup \\
    \underline{\omega}_{M_0} & \xhookrightarrow{\ \mathrm{can}\ } & \operatorname{Aut}(M_0) \simeq \mathfrak{S}_J \\
    \| & & \\
    \underline{\omega}_J \simeq \mathfrak{S}_J^+ \setminus \{\mathrm{id}_J\} & &
  \end{array}
\]
\emph{NB} \struck{On a} si $\rho \in \underline{\omega}$, $\rho' \overset{\mathrm{def}}{=} \rho^{-1}$, $\rho + \rho'$ \struck{$+$} $+\,\mathrm{id} = 0$ i.e. $\rho + \rho' = \mathrm{id}_{M_0}$, i.e. l'hom. $\mathbb{F}_2^{\underline{\omega}} \hookrightarrow \operatorname{End}(M_0)$ \uncertain{applique} $\mathbb{1}$ sur $1$.

Si $i, j \in J$, $i \neq j$, alors $\sigma_i, \sigma_j, \rho, \rho' \in \operatorname{End}(M_0)$ sont lin. indépendants sur $\mathbb{F}_2$ donc forment une base, \struck{et} \ill{} que \struck{$x_i, x_j$} en forment une de $M_0$, [et $\rho, \rho'$] de $\mathbb{F}_2^{\underline{\omega}}$\ldots
\note{les quatre dernières lignes, d'une écriture plus serrée, sont au bas à droite de la page ; la fin en est d'une lecture incertaine.}

\page{27}
On a
\[
  (15) \qquad
  \left\{
  \begin{array}{l}
    M_0 \text{ et } \mathbb{F}_2^{\underline{\omega}} \text{ sont orthogonaux l'un de l'autre} \\
    \text{dans } \underline{\operatorname{End}}(M_0) \text{ pour la forme } \operatorname{Tr}(uv)
  \end{array}
  \right.
\]
\[
  (16) \qquad \mathfrak{sl}(M_0) = M_0 \oplus \mathbb{F}_2\, \mathrm{id}_{M_0}
\]
\[
  (17) \qquad \mathbb{F}_2^{\underline{\omega}} \cap \mathfrak{sl}(M_0) \simeq \mathbb{F}_2 \cdot \mathrm{id}_{M_0} \qquad \ldots \qquad ]
\]
\note{sous le symbole de (16), un trait le relie à « $\operatorname{Ker}(\operatorname{End}(M_0) \xrightarrow{\mathrm{tr}} \mathbb{F}_2)$ », qui en est la définition ; le symbole, écrit sur un autre noirci, a la forme d'un « $\gamma\ell$ » cursif. Il est rendu $\mathfrak{sl}$ d'après cette définition, et parce que la page~31 l'oppose à un « $\mathfrak{gl}(M_0)$ » tracé autrement. Le crochet fermant clôt la parenthèse ouverte page~26 par « [Notons\ldots ».}

Ainsi, la donnée d'un torseur $P_M$ sous $V_J(M_0) \simeq \operatorname{End}(M_0)$, équivaut à celle d'un couple de torseurs, \struck{\ill{}} l'un sous $M_0$, l'autre sous $\mathbb{F}_2^{\underline{\omega}}$, soient $I_M$ et $S_M$.

\emph{Théorème} Le foncteur $M \mapsto (M_0, I_M, S_M)$
\[
  \bigl(\mathbb{Z}/4\mathbb{Z}\text{-modules libres de rang } 2\bigr)
  \longrightarrow
  \text{triplets } (M_0, I, S) \ \text{avec}\
  \left\{
  \begin{array}{l}
    M_0 \text{ vect. de rg } 2 \text{ sur } \mathbb{F}_2 \\
    I \ \ M_0\text{-torseur} \\
    S \ \ \text{torseur sous } \mathbb{F}_2^{\underline{\omega}(M_0)}
  \end{array}
  \right.
\]
est une équivalence de groupoïdes. (18)
\note{le numéro (18), encerclé et surchargé, est dans la marge, en face de la flèche ; sous « triplets », un premier libellé « où $M_0$ un vect. de rg 2 » est biffé et remplacé par « avec » et l'accolade ; dans le triplet, les indices de $I$ et $S$ sont noircis.}

\emph{Cor.} \struck{\ill{}} On a un isom. de suites exactes canoniques ($\forall M$ libre de rg $2$ sur $\Lambda = \mathbb{Z}/4\mathbb{Z}$)

\begin{tikzcd}[column sep=small]
    1 \arrow[r] & \operatorname{End}_{\mathbb{F}_2}(M_0) \arrow[r] & \operatorname{GL}_\Lambda(M) \arrow[r] & \operatorname{GL}(M_0) \simeq \mathfrak{S}_J \arrow[r] & 1 \\
    1 \arrow[r] & M_0 \times \mathbb{F}_2^{\underline{\omega}} \arrow[r] \arrow[u, "\wr"'] & \operatorname{Aut}(M_0, I_M, S_M) \arrow[r] \arrow[u, "\wr"'] & \operatorname{GL}(M_0) \arrow[r] \arrow[u, no head, "\|" description] & 1
  \end{tikzcd}

Ici l'hom. $\operatorname{End}(M_0) \to \operatorname{GL}(M)$ est donné par la condition que le composé avec l'hom. $\operatorname{End}_\Lambda(M) \to \operatorname{End}_{\mathbb{F}_2}(M_0)$ soit donné par $u \mapsto \mathrm{id}_M + 2u$.
\note{la phrase finale est transcrite telle quelle ; le sens attendu est que $u \in \operatorname{End}(M_0)$ s'envoie sur $\mathrm{id}_M + 2\tilde u$, $\tilde u$ un relèvement de $u$ à $\operatorname{End}_\Lambda(M)$.}

\marginal{Équivaut \uncertain{en outre} à l'isom. $\operatorname{GL}(2, \mathbb{Z}/4\mathbb{Z}) \simeq \operatorname{GL}(2, \mathbb{F}_2) \cdot M_2(\mathbb{F}_2)$ (produit semi-direct), plus l'isomorphisme de $\mathfrak{S}_3$-modules $M_2(\mathbb{F}_2) \simeq \mathbb{F}_2^2 \times \mathbb{F}_2^{\underline{\omega}_0}$ ($\mathfrak{S}_3 \simeq \operatorname{GL}(2, \mathbb{F}_2)$), où \ill{} \ill{} \ill{} \ill{} \ill{}\ldots\ On a \struck{\ill{}} $\operatorname{SL}(2, \mathbb{Z}/4\mathbb{Z}) \simeq \{ g \cdot u \in \mathfrak{S}_3 \cdot M_2(\mathbb{F}_2) \mid \operatorname{sg}(g)\, e(\operatorname{Tr} u) = 1 \}$ $(19)$ où $e : \mathbb{F}_2 \xrightarrow{\;\sim\;} \{\pm 1\}$.}
\note{la note marginale est écrite obliquement dans la marge gauche, du milieu au bas de la page, et une partie (autour de $\operatorname{SL}(2, \mathbb{Z}/4\mathbb{Z})$) est chargée de ratures ; les formules sont plus sûres que les mots. Le numéro (19) est de l'auteur.}

\page{28}\note{p.~14 de l'auteur.}
La donnée de $(M_0, I)$ équivaut, par le §~1, à la donnée d'un ens. à quatre éléments $I$, on aura $M_0 \simeq V_I$, $M_0^* \simeq \mathfrak{P}_{2,2}(I)$. \struck{Par ailleurs} \add{D'autre} part, \struck{la donnée d'un} \ldots\ un ens. à deux éléments $C$, et d'un torseur $S$ sous $\mathbb{F}_2^C$, équivaut à la donnée d'un cube $Q$, $S$ étant l'ens. des sommets de $Q$, $C$ l'ens. de ses codiagonales. Donc le \struck{catégorie} \add{groupoïde} du deuxième membre de $(18)$, est équivalent à celui des triples $(I, Q; k)$, où $I$ est un ens. à quatre éléments, $Q$ un cube, et $k$ un isom.
\[
  (20) \qquad k : \underset{\substack{\text{ens. des}\\ \text{orientations}\\ \text{de } I}}{\underline{\omega}(I)} \simeq \underset{\substack{\text{ens. des}\\ \text{codiagonales}\\ \text{de } Q}}{C(Q)} .
\]
\note{le passage biffé à la quatrième ligne (« la donnée d'un\ldots ») est souligné puis barré ; le « \ldots » marque le raccord.}

On trouve donc

\emph{Corollaire 2} On a une équivalence canonique de groupoïdes
\[
  (21) \qquad
  \Bigl( \underset{\mathbb{Z}/4\mathbb{Z}}{\underline{\Lambda}}\text{-modules libres de rg } 2 \Bigr)
  \xrightarrow{\;\approx\;}
  \text{triplets } (I, Q, k), \ \text{où}\
  \left\{
  \begin{array}{l}
    I \ \text{ens. à 4 éléments} \\
    Q \ \text{cube combinatoire} \\
    k : \underline{\omega}(I) \simeq C(Q)
  \end{array}
  \right.
\]

\emph{Corollaire 3} Si $M$ est associé à $I$, $Q$, on a un diagramme cartésien canonique
(22)


\begin{tikzcd}[column sep=small]
    & \operatorname{GL}_\Lambda(M) \arrow[dl] \arrow[dr] & \\
    \mathfrak{S}_I \arrow[dr, "\mathrm{sg}"'] & & \mathfrak{D}_Q \arrow[dl, "\chi_C"] \\
    & \{\pm 1\} &
  \end{tikzcd}

\note{à côté de $\chi_C$, l'auteur précise : « (opération sur l'ens. des codiagonales) ».}

\page{29}
s'inscrivant dans le diagramme de suites exactes (\uncertain{auquel} il donne naissance)
(23)


\begin{tikzcd}[column sep=small, row sep=small, nodes={font=\scriptsize}]
    & & & 1 \arrow[d] & & & \\
    & 1 \arrow[dr] & & \mathfrak{S}_I^+ \times \mathbb{F}_2^{\underline{\omega}} \arrow[dl, "\mathrm{pr}_1"'] \arrow[dr, "\mathrm{pr}_2"] \arrow[dd] & & 1 \arrow[dl] & \\
    & & \mathfrak{S}_I^+ \arrow[dl, Rightarrow] \arrow[dr] & & \mathbb{F}_2^{\underline{\omega}} \arrow[dl] \arrow[dr, Rightarrow] & & 1 \arrow[dl] \\
    1 \arrow[r] & \mathfrak{S}_I^+ \arrow[dr] & & \operatorname{GL}_\Lambda(M) \arrow[dl] \arrow[dr] \arrow[dd, "\text{cf.\ (24)}"] & & \mathbb{F}_2^{\underline{\omega}} \simeq V_{S_Q} \arrow[dl] & \\
    & & \mathfrak{S}_I \arrow[dl] \arrow[dr, "\mathrm{sg}"'] & & \mathfrak{D}_Q \arrow[dl, "\chi_C"] \arrow[dr] & & \\
    & 1 & & \{\pm 1\} \arrow[dl] \arrow[d] \arrow[dr] & & 1 & \\
    & & 1 & 1 & 1 & &
  \end{tikzcd}

\note{diagramme redessiné. Les deux flèches doubles (de $\mathfrak{S}_I^+$ vers la seconde copie de $\mathfrak{S}_I^+$, et de $\mathbb{F}_2^{\underline{\omega}}$ vers $\mathbb{F}_2^{\underline{\omega}} \simeq V_{S_Q}$) sont tracées comme telles sur la page et marquent des identifications. L'étiquette de la flèche verticale de $\operatorname{GL}_\Lambda(M)$ vers $\{\pm 1\}$ est peu lisible ; la lecture « cf.\ (24) » (renvoi à la formule qui suit) s'appuie sur le diagramme (25) de la page~30, où elle est plus nette.}

\emph{Proposition} Le \add{noyau ($\simeq \mathfrak{S}_I^+$) de l'} hom. canonique $\operatorname{GL}_\Lambda(M) \to \mathfrak{D}_Q$ \struck{induit un iso} est aussi égal au groupe des commutateurs du sous-groupe \add{distingué} $\operatorname{SL}_\Lambda(M)$ d'indice $2$ de $\operatorname{GL}_\Lambda(M)$ (NB $\Lambda^* = (\mathbb{Z}/4\mathbb{Z})^* = \{\pm 1\}$), qui est \add{dans} d'indice $4$ dans $\operatorname{SL}_\Lambda(M)$ --- le quotient $(\operatorname{SL}_\Lambda(M))_{\mathrm{ab}}$ est cyclique d'ordre $4$ \struck{isomorphe à $\mathbb{Z}/4\mathbb{Z}$}, et l'hom.
\[
  \operatorname{GL}_\Lambda(M) \xrightarrow{\;\det\;} \{\pm 1\}
\]
est le composé \struck{\ill{}}
\[
  (24) \qquad \operatorname{GL}_\Lambda(M) \longrightarrow \mathfrak{D}_Q \underset{\substack{\text{noyau } \mathfrak{D}_Q^+,\\ \text{cyclique d'ordre 4}}}{\underbrace{\xrightarrow{\;\det\;}}} \{\pm 1\} .
\]
\note{l'étiquette « det » de la seconde flèche de (24) est écrite sur un mot noirci.}

\emph{Corollaire \ill{}} On a un sous-diagramme de suites exactes de $(23)$, obtenu ainsi :

\page{30}\note{p.~15 de l'auteur, chiffrée « (15) » en tête de page.}
(25)


\begin{tikzcd}[column sep=small, row sep=small, nodes={font=\scriptsize}]
    & & & 1 \arrow[d] & & & \\
    & 1 \arrow[dr] & & \mathfrak{S}_I^+ \times \{1, -\mathrm{id}_M\} \arrow[dl, "\mathrm{pr}_1"'] \arrow[dr, "\mathrm{pr}_2"] \arrow[dd] & & 1 \arrow[dl] & \\
    & & \mathfrak{S}_I^+ \arrow[dl, Rightarrow] \arrow[dr] & & \{1, -\mathrm{id}_M\} \arrow[dl] \arrow[dr, "\wr"] & & 1 \arrow[dl] \\
    1 \arrow[r] & \mathfrak{S}_I^+ \arrow[dr] & & \operatorname{SL}_\Lambda(M) \arrow[dl] \arrow[dr] \arrow[dd, "\text{cf.\ (24)}"] & & \{1, \underline{a}_Q\} \simeq \mathbb{F}_2 \arrow[dl] & \\
    & & \mathfrak{S}_I \arrow[dl] \arrow[dr, "\mathrm{sg}"'] & & \mathfrak{D}_Q^+ \arrow[dl] \arrow[dr] & & \\
    & 1 & & \{\pm 1\} \arrow[dl] \arrow[d] \arrow[dr] & & 1 & \\
    & & 1 & 1 & 1 & &
  \end{tikzcd}

\note{diagramme redessiné, de même forme que (23). En haut, deux flèches qui aboutissaient à $\mathfrak{S}_I^+ \times \{1, -\mathrm{id}_M\}$ de part et d'autre de la flèche verticale sont biffées de traits serrés et ne sont pas reprises. L'exposant de $\mathfrak{S}_I$ dans le terme du haut est surchargé ; il est lu $+$. Le « $1$ » qui aboutit à $\{1, \underline{a}_Q\}$ est écrit sur un autre signe, et le premier élément de cet ensemble sur un symbole noirci.}

\emph{Corollaire 2} L'hom. can. $\operatorname{GL}(M) \to \mathfrak{D}_Q$ induit un isomorphisme
\[
  (26) \qquad \operatorname{GL}_\Lambda(M)_{\mathrm{ab}} = \operatorname{GL}_\Lambda(M) / \mathfrak{S}_I^+ \times \underset{\text{centre de } \operatorname{GL}_\Lambda(M)}{\underbrace{\{1, -\mathrm{id}_M\}}} \xrightarrow{\;\sim\;} \underset{\substack{\text{vectoriel de rang 2 sur } \mathbb{F}_2\\ \text{et même \underline{canon.} isom. à}\\ \mathbb{F}_2 \times \mathbb{F}_2}}{(\mathfrak{D}_Q)_{\mathrm{ab}}}
\]
\marginal{\emph{NB.} $\left\{ \begin{array}{l} \det g = \chi_{\underline{\omega}}(\dot g) \\ \mathrm{sg}\, g = \chi_C(\dot g) \\ \det g \cdot \mathrm{sg}\, g = \chi_D(\dot g) \end{array} \right.$ pour $g \in \operatorname{GL}(M)$, $\dot g$ im. can. $\in \mathfrak{D}_Q$}
\note{le numéro (26) est à demi effacé. Dans la note marginale, l'indice du premier $\chi$ est écrit sur un signe noirci ; la lecture $\chi_{\underline{\omega}}$ est incertaine.}

On peut réunir \add{l'essentiel de l'information des} (23) (25) dans un diagramme de sous-groupes de $\operatorname{GL}_\Lambda(M)$, où on a posé pour abréger
\[
  (27) \qquad \mathfrak{z} = \{1, -\mathrm{id}_M\} \ \bigl(= \operatorname{Cent}(\operatorname{GL}_\Lambda(M)) = \operatorname{Cent}(\operatorname{SL}_\Lambda(M))\bigr)
\]
\note{la lettre qui désigne le centre est un $\mathfrak{z}$ cursif (« ʒ »).}

\page{31}
(28)


\begin{tikzcd}[column sep=small, row sep=small, nodes={font=\scriptsize}]
    & & & \operatorname{GL}_\Lambda(M) \simeq [\mathfrak{S}_I \cdot \mathbb{F}_2^{\underline{\omega}}] & & \\
    & & \mathfrak{S}_I^+ \times \mathbb{F}_2^{\underline{\omega}} \ [\mathfrak{S}_I \times \mathfrak{z}] \arrow[ur, "2"] & & \operatorname{SL}_\Lambda(M) \arrow[ul, "2"'] & \\
    & M_0 \times \mathbb{F}_2^{\underline{\omega}} \simeq \mathfrak{gl}(M_0) \arrow[ur, "3"] & & \mathfrak{S}_I^+ \times \mathfrak{z} \arrow[ul, "2"'] \arrow[ur, "2"] \arrow[uu, dashed] & & \\
    \mathbb{F}_2^{\underline{\omega}} \arrow[ur, "4"] & & M_0 \times \mathfrak{z} \simeq \mathfrak{sl}(M_0) \arrow[ul, "2"'] \arrow[ur, "3"] & & \mathfrak{S}_I^+ \arrow[ul, "2"'] & \\
    & \mathfrak{z} \arrow[ul, "2"'] \arrow[ur, "4"] & & M_0 \arrow[ul, "2"'] \arrow[ur, "3"] & & \\
    & & 1 \arrow[ul, "2"'] \arrow[ur, "4"] & & &
  \end{tikzcd}

\note{diagramme redessiné : le treillis des sous-groupes de $\operatorname{GL}_\Lambda(M)$, chaque flèche (une inclusion) portant l'indice. La flèche pointillée monte de $\mathfrak{S}_I^+ \times \mathfrak{z}$ à $\operatorname{GL}_\Lambda(M)$ en passant par le crochet $[\mathfrak{S}_I \times \mathfrak{z}]$ ; les crochets semblent indiquer une autre description du même groupe. Des accolades, qui ne sont pas reproduites, désignent des quotients : l'une, à gauche, embrasse la bande supérieure du treillis et porte $\mathfrak{S}_J \simeq \operatorname{GL}_{\mathbb{F}_2}(M_0)$ (avec un second $\mathfrak{S}_J$ au-dessus) ; à droite, une accolade entre $\operatorname{SL}_\Lambda(M)$ et $\mathfrak{S}_I^+$ porte $\mathfrak{D}_Q^+ \simeq (\operatorname{SL}_\Lambda(M))_{\mathrm{ab}} \simeq \bigwedge^2_\Lambda M$ (suivi d'un mot noirci), une grande accolade $\mathfrak{D}_Q$ embrasse tout le côté droit, et une accolade sous $\mathfrak{S}_I^+$ porte $\mathfrak{S}_J^+$. $\mathfrak{gl}(M_0)$ et $\mathfrak{sl}(M_0)$ rendent deux symboles cursifs que la page distingue (cf.\ la note de la page~27).}

\emph{NB}
\[
  \begin{array}{l}
    \underline{\omega} \subset \mathbb{F}_2^{\underline{\omega}} = \mathbb{F}_2^C = V_S \subset \mathfrak{D}_Q \quad (C_Q \simeq \underline{\omega}) \\
    D_Q \simeq \underline{\omega} \wedge \mu \subset \mathbb{F}_2^D = V_A \subset \mathfrak{D}_Q \\
    \underline{\omega}_Q \simeq \mu \subset \mathfrak{D}_Q^+ \subset \mathfrak{D}_Q
  \end{array}
  \qquad
  \begin{array}{l}
    \underline{\omega} \simeq V_S \setminus \mathfrak{z}_Q \longleftrightarrow \mathrm{sg}(g) \\
    \underline{\omega} \wedge \mu \simeq V_A \setminus \mathfrak{z}_Q \longleftrightarrow \mathrm{sg}(g)\det(g) \\
    \mu \simeq \mathfrak{D}_Q^+ \setminus \mathfrak{z}_Q \longleftrightarrow \det(g)
  \end{array}
\]
où
\[
  \mu \simeq \bigl(\textstyle\bigwedge^2 M\bigr)^* = \operatorname{Or}(M), \qquad \mathfrak{z}_Q = \{1, \underline{a}_Q\} = \operatorname{Cent} \mathfrak{D}_Q
\]
\marginal{$\operatorname{GL}_\Lambda(M)' \simeq \mathfrak{S}_I \times \mathbb{F}_2 \simeq \operatorname{Aut}(\Delta)$, où $\Delta$ est l'octaèdre non orienté ; ensemble de sommets $\Delta = M^{\#}/\{\pm 1\}$\ldots}
\note{ces formules, avec la note en haut à droite, sont écrites au-dessus du diagramme et en sont séparées par un trait. Le « $C_Q$ » est écrit au-dessus de $\underline{\omega}$ avec un signe $\simeq$ vertical. Dans la note de droite, « octaèdre non orienté » est d'une lecture incertaine, et le symbole de $M^{\#}$ (les éléments de $M$ hors de $2M$, cf.\ (29) page~32) n'est pas sûr. Dans $\mathrm{sg}(g)\det(g)$, un premier « sg » est surchargé.}

\emph{NB} Dans ce diagramme, toutes les flèches sont injectives et les cinq carrés qui y apparaissent, donc aussi les composés, sont cartésiens. On a marqué sur chaque flèche l'indice \struck{\ill{}} de l'inclusion des sous-groupes qu'elles représentent.

Les iso.
\[
  \begin{aligned}
    \operatorname{GL}_\Lambda(M) &\simeq \mathfrak{S}_I \cdot \mathbb{F}_2^{\underline{\omega}} && (\tfrac12\text{ direct}) \ \text{et} \\
    \operatorname{Ker}(\mathrm{sg} \cdot \det) &\simeq \mathfrak{S}_I \times \mathbb{F}_2 && (\text{prod. direct})
  \end{aligned}
\]
sont valables quand $M$ est « $\tfrac12$déployé » i.e. le torseur $P_M$ est trivialisé, dans ce cas
\[
  \begin{aligned}
  \operatorname{GL}_\Lambda(M) &\simeq \underset{\mathfrak{S}_J}{\underbrace{\operatorname{GL}_{\mathbb{F}_2}(M_0)}} \cdot (M_0 \times \mathbb{F}_2^{\underline{\omega}}) \\
  &\simeq \bigl(\underset{\mathfrak{S}_I}{\underbrace{\operatorname{GL}_{\mathbb{F}_2}(M_0) \cdot M_0}}\bigr) \times \mathbb{F}_2^{\underline{\omega}}
  \end{aligned}
\]
mais ces isom. dépendent du choix d'une trivialisation de $P_M$\ldots
\note{« $\tfrac12$ direct » et « $\tfrac12$déployé » : le signe $\tfrac12$ vaut « semi- » (produit semi-direct, semi-déployé). Au-dessus de la dernière formule, un « $=$ » isolé.}

\section{6. Relations avec l'octaèdre non orienté $\Delta_M$.}

\page{32}\note{p.~16 de l'auteur.}
Soit
\[
  \begin{aligned}
  (29) \qquad M^{\#} &= \{ x \in M \mid x^0 \ (\text{image dans } M_0) \neq 0 \} \\
  &= \text{image inverse de } J = M_0^{\#} \subset M_0 \text{ par } M \to M_0
  \end{aligned}
\]
\struck{donc on} donc $M^{\#}$ est un torseur relatif sur $M_0^{\#} = J$, de groupe $(M_0)_J$ (i.e. famille constante de groupes isom. à $M_0$, indexée par $J$). Le groupe $(\mathbb{Z}/4\mathbb{Z})^* = \{\pm 1\}$ opère sur $M^{\#}$, de façon compatible avec la projection $M^{\#} \to J$, on pose
\[
  (30) \qquad \Delta_M = M^{\#}/(\mathbb{Z}/4\mathbb{Z})^* = M^{\#}/\{\pm 1\} = \mathbb{P}(M)(\mathbb{Z}/4\mathbb{Z})
\]
ce sont les points \uncertain{à valeurs} dans $\Lambda = \mathbb{Z}/4\mathbb{Z}$ de la droite projective sur $\Lambda$ définie par $M$, et l'application \add{can.} \struck{\ill{}}
\[
  (31) \qquad \Delta_M \longrightarrow J
\]
s'identifie à l'application de restriction, ou réduction mod~$2$
\[
  \mathbb{P}(M)(\mathbb{Z}/4\mathbb{Z}) \longrightarrow \mathbb{P}(M)(\mathbb{F}_2) .
\]
C'est une application de degré $2$, qui définit donc sur $\Delta$ la structure de (l'ens. des sommets) d'un octaèdre, dont $J$ est l'ens. des diagonales. Les $12$ arêtes de l'octaèdre sont les éléments de \struck{le groupe $\operatorname{GL}(M)$ opère sur} $\mathfrak{P}_2(\Delta)$ (qui a $\frac{6 \cdot 5}{2} = 15$ éléments) qui ne sont pas \struck{des diagonales} l'une des trois diagonales, i.e. des trois fibres de $(31)$. Les $8$ faces $f$ de l'octaèdre sont les parties à trois éléments de $\Delta$ telles que l'application induite par $(31)$ soit une bijection
\[
  (32) \qquad f \xrightarrow{\;\sim\;} J \ ;
\]
elles s'identifient aux sections de $(31)$, tandis que les arêtes s'identifient aux sections partielles sur des parties de $J$ à deux éléments, et \uncertain{Ce\ldots}
\note{le signe de $M^{\#}$ est un astérisque ou un dièse ; il est rendu $M^{\#}$ pour le distinguer du groupe des unités $(\mathbb{Z}/4\mathbb{Z})^*$ de la même page. « à valeurs » est une lecture incertaine. Au-dessus de « faces », un « $f$ » et un signe noirci sont ajoutés ; au-dessus de « trois éléments », un « $6$ » (?). La phrase se poursuit page~33.}
\page{33}
\uncertain{ment}) les sommets s'identifient aux sections partielles sur des parties \add{de $J$} réduites à un élément --- la relation d'incidence est celle de prolongement des sections, ou encore l'inclusion des parties de $\Delta$.

Le groupe $G = \operatorname{GL}(M)$ opère sur $\Delta$ par \uncertain{automorphismes} d'octaèdre, on voit que $\mathfrak{z} = \{\pm 1\}$ opère trivialement, \struck{\ill{}} et \uncertain{ce} que c'est là le noyau de l'opération, on a donc un hom. injectif
\[
  \underset{\substack{\wr\\ \operatorname{GL}(M)/\{\pm 1\}}}{\operatorname{GL}(M)'} \longrightarrow \operatorname{Aut}_{\mathrm{oct}}(\Delta)
\]
\marginal{\emph{NB} $\operatorname{GL}(M)' \xrightarrow{\;\sim\;} \operatorname{GP}(M)$ (groupe projectif de \uncertain{rang}~$1$)}
qui est un isomorphisme, car les deux groupes sont de même ordre $24 \times 2 = 48$, donc on trouve
\[
  (33) \qquad
  \underset{\substack{\|\\ \operatorname{GL}(M)/(\mathfrak{z} = \{\pm 1\})\\ \wr\\ \operatorname{GP}(M)}}{\operatorname{GL}(M)'} \xrightarrow{\;\sim\;} \operatorname{Aut}_{\mathrm{oct}}(\Delta) \simeq \operatorname{Aut}^+(\Delta) \times \underset{\substack{\text{antipodisme}\\ \text{centre de } \operatorname{Aut}_{\mathrm{oct}}(\Delta)}}{\{1, \underline{a}_\Delta\}}
\]
le centre de $\operatorname{GL}(M)'$ \struck{correspondant} \add{n'est autre} d'ailleurs que le sous-groupe $\mathbb{F}_2^{\underline{\omega}}/\mathbb{F}_2 \cdot \mathrm{id}_{M_0}$ de $\operatorname{GL}(M)'$ (car le centre de $\operatorname{GL}(M_0) \simeq \mathfrak{S}_J$ est nul, donc le centre de $\operatorname{GL}(M)'$ est contenu dans $\underline{\operatorname{End}}(M_0)/\mathfrak{z} = \{1_{M_0}, \ldots\} \simeq M_0 \times \mathbb{F}_2^{\underline{\omega}}/\mathbb{F}_2$, et s'identifie au sous-groupe des invariants sous $\mathfrak{S}_J \simeq \operatorname{GL}_{\mathbb{F}_2}(M_0)$ de ce groupe, c'est $\mathbb{F}_2^{\underline{\omega}}/\mathbb{F}_2$.
\note{le premier mot de la page achève le mot coupé en bas de la page~32. Dans « $\underline{\operatorname{End}}(M_0)/\mathfrak{z} = \{\ldots\}$ », le second élément de l'ensemble est surchargé et n'est pas lu. La parenthèse ouverte par « (car » n'est pas refermée.}

La donnée d'un octaèdre $\Delta$ ayant $J$ comme ens. des diagonales équivaut à celle d'un torseur
\page{34}\note{p.~17 de l'auteur.}
sous $\mathbb{F}_2^J$, que je vais désigner par $P_{M,\Delta}$
\[
  (34) \qquad \underset{\text{ens. des faces de } \Delta}{\underbrace{F(\Delta_M)}} = P_{M,\Delta} \quad \text{torseur sous } \mathbb{F}_2^J \qquad (J = M_0^{\#})
\]
Or on a
\[
  (35) \qquad \mathbb{F}_2^J \simeq \underset{\substack{\wr\\ M_0}}{V_J(\mathbb{F}_2)} \times \underset{\text{diag}}{\mathbb{F}_2}
\]
et l'opération sur $P_{M,\Delta}$ de l'élément non nul \struck{\ill{}} du facteur $\mathbb{F}_2$ n'est autre que l'antipodisme. Ceci posé on a un isom. canonique
\[
  (36) \qquad P_{M,\Delta} = F(\Delta_M) \xleftarrow{\;\sim\;} P_M/\mathfrak{z} \qquad \underset{\text{centre de } \operatorname{GL}(M)}{\mathfrak{z} = \{\pm \mathrm{id}_M\}}
\]
provenant de l'homomorphisme évident
\[
  (37) \qquad
  \begin{array}{l}
    P_M = \{ (\xi_i)_{i \in J} \in M^{\# J} \mid \xi_i^0 = i_* \ \forall i \in J,\ \sum \xi_i = 0 \} \\
    \quad \big\downarrow \qquad (\xi_i) \longmapsto \{(\dot\xi_i)\} \quad \text{où } \dot\xi = \Lambda^* \xi = \pm \xi \in \Delta \text{ pour } \xi \in M^{\#} \\
    F(\Delta_M) = \{ f \in \mathfrak{P}_3(\Delta_M) \mid \underset{\text{induit par } \Delta \to J}{\underbrace{f \to J}} \text{ bijectif} \}
  \end{array}
\]
Ceci nous montre à nouveau que $\operatorname{Aut}(\Delta_M) \simeq \operatorname{GL}(M)'$, de façon précise on a l'iso. de suites exactes
(38)


\begin{tikzcd}[column sep=small]
    1 \arrow[r] & \operatorname{End}(M_0)/\mathbb{F}_2 \cdot \mathrm{id}_{M_0} \arrow[r] \arrow[d, "\wr"] & \operatorname{GL}(M)' \arrow[r] \arrow[d, "\wr"] & \operatorname{GL}_{\mathbb{F}_2}(M_0) \simeq \mathfrak{S}_J \arrow[r] \arrow[d, "\wr"] & 1 \\
    1 \arrow[r] & \mathbb{F}_2^J \simeq M_0 \times \mathbb{F}_2 \arrow[r] & \operatorname{Aut}_{\mathrm{oct}}(\Delta_M) \arrow[r] & \mathfrak{S}_J \arrow[r] & 1
  \end{tikzcd}

\note{au-dessus de $\operatorname{End}(M_0)/\mathbb{F}_2 \cdot \mathrm{id}_{M_0}$, une accolade porte « $M_0 \times \mathbb{F}_2$ » ; la flèche verticale de gauche porte en outre un astérisque, sans renvoi visible.}

\ill{} \ill{} \ill{} \ill{} associés à $(\Delta, \underline{a}_\Delta)$ ou $(\Delta \to J)$ sont \ill{} \ill{} trois ens. à deux éléments remarquables, \ill{} aux trois caractères \add{non triviaux} de $\operatorname{Aut}_{\mathrm{oct}}(\Delta)$ \ill{} trois hom. non nuls de \struck{$\operatorname{Aut}_{\mathrm{oct}}(\Delta)$}
\[
  \operatorname{Aut}_{\mathrm{oct}}(\Delta)_{\mathrm{ab}} \simeq \bigl(\operatorname{Aut}^+_{\mathrm{oct}}(\Delta) \times \{1, \underline{a}_\Delta\}\bigr)_{\mathrm{ab}} \simeq \mathbb{F}_2 \times \mathbb{F}_2
\]
sont respectivement les ensembles
\marginal{\emph{NB} La donnée d'un octaèdre (pas orienté) $\Delta$ équivaut à celle d'un couple $(I, \ldots)$, où $I$ est un ens. à quatre éléments \ill{} \ill{} \ill{} \ill{} \ill{} \ill{} octaèdre \ill{} \ill{} \ill{} \ill{} de ses \ill{} \ill{} \ill{} trois \ill{} \ill{} \ill{} \ill{} \ill{} (l'un des \ill{}, c'est l'oct. \emph{orienté} \ill{} \ill{}) $= \underline{\omega}(\Delta)$ \ill{} deux éléments \ill{} i.e. \ill{} \ill{} $I$ \ill{} (\ill{} $(39)$ \ill{}) que \ill{} \ill{} \ill{}}
\note{la note marginale, écrite obliquement dans le coin inférieur gauche, déborde sur le début des lignes du corps qu'elle recouvre en partie ; le début de la phrase du corps (« \ldots associés à $(\Delta, \underline{a}_\Delta)$ ») est ainsi perdu. La note n'est lue que par fragments ; on y distingue « $\underline{\omega}_I$ » et un renvoi à (39). Au-dessus de « caractères », « non triviaux » est ajouté en interligne.}
\page{35}
\[
  (40) \qquad
  \left\{
  \begin{array}{l}
    C_\Delta \ (\text{« couleurs »}) : \text{ ens.\ des 2 tétraèdres (échangés par l'antipodisme)} \\
    \quad \text{inscrits dans le cube qui correspond à } \Delta, \\
    \quad \text{ou encore des deux couleurs (échiquier\ldots) des faces de l'octaèdre } \Delta \\
    \quad \text{correspond au car.\ \ldots\ de } \operatorname{Aut}^+_{\mathrm{oct}}(\Delta) \times \{1, \underline{a}_\Delta\} \\[4pt]
    \operatorname{Or}(\Delta) : \text{ ens.\ des deux orientations de l'octaèdre } \Delta \\
    \quad \simeq \text{ correspond au car.\ non trivial de } \{1, \underline{a}_\Delta\} \\[4pt]
    \underline{\omega}_\Delta = \underline{\omega}_{I_\Delta} : \text{ ens.\ des deux orientations de } I_\Delta,
      \text{ ou encore de } J_\Delta \\
    \quad \text{correspond au caractère non trivial de } \operatorname{Aut}^+_{\mathrm{oct}}(\Delta)
  \end{array}
  \right.
\]
\emph{NB} On a $F(\Delta) \simeq C_\Delta \times I_\Delta$, où $I_\Delta = F(\Delta)/\underline{a}_\Delta$ est l'ens. des quatre directions de faces
\note{le tableau (40) est redessiné en colonnes. Dans la première colonne, deux figures : un cube dont les huit sommets sont marqués alternativement de points pleins et de cercles (les deux tétraèdres inscrits), et, dessous, un octaèdre dont les faces sont hachurées une sur deux (les deux couleurs, en damier). Le symbole $C_\Delta$ est écrit sur un signe noirci ; en tête de la troisième colonne, un mot est biffé avant $\underline{\omega}_\Delta$. Dans la première colonne, le « \ldots » devant « de $\operatorname{Aut}^+_{\mathrm{oct}}$ » tient la place d'un mot biffé, illisible.}

qui satisfont à un iso. canonique
\[
  (41) \qquad \operatorname{Or}(\Delta) \wedge \underline{\omega}_\Delta \wedge C_\Delta \simeq \mathbb{1}_{\mathbb{F}_2} .
\]
On a un isom.
\[
  (42) \qquad
  \left\{
  \begin{array}{l}
    C_\Delta = \bigwedge_{i \in J} \Delta_i \quad \text{quotient de } \overbrace{\prod_{i \in J} \Delta_i}^{\text{torseur sous } \mathbb{F}_2^J} = F(\Delta) \quad \text{par } \underset{\substack{\wr\\ \underbrace{V_J(\mathbb{F}_2)}_{M_0} \times \mathbb{F}_2}}{\mathbb{F}_2^J} \xrightarrow{\text{(somme)}} \mathbb{F}_2 \\[4pt]
    \text{i.e. } C_\Delta = \underbrace{F(\Delta)/M_0}_{(P_M/(M_0 \times \mathfrak{z}) \text{ dans le cas qui nous occupe})}
  \end{array}
  \right.
\]
correspond au car. $\mathrm{sg}(g) \cdot \det(g)$ dans le cas où $\Delta$ est défini par un $M$)
\marginal{donc dans ce cas on a $C_\Delta \simeq \mu \wedge \underline{\omega}$}
\note{après $P_M/(M_0 \times \mathfrak{z})$, un mot est noirci. La parenthèse qui se ferme après « un $M$ » n'a pas d'ouvrante visible.}
\[
  (43) \qquad \underline{\omega}_\Delta = \underline{\omega}_{I_\Delta} = \underline{\omega}_J \qquad \text{l'ens. à deux éléments associé à } J
\]
(correspond au caractère $\mathrm{sg}(g)$ dans le cas d'un $M$)
\[
  (44) \qquad \operatorname{Or}(\Delta) \simeq \underset{\substack{\underline{\omega}_J\\ \Delta^+}}{\underbrace{\underline{\omega}_\Delta}} \wedge C_\Delta \qquad \left(\begin{array}{l}\text{correspond au caractère } \det(g) \text{ dans le cas d'un } M) \\ \text{d'où } \boxed{\operatorname{Or}(\Delta) \simeq \mu \overset{\mathrm{def}}{=} (\det M)^*}\end{array}\right.
\]
\note{devant $\operatorname{Or}(\Delta)$ en (44), un signe noirci ; un autre devant le numéro (43). Le $\Delta^+$ sous l'accolade est relié par un trait à « orienté » de la ligne suivante.}

L'octaèdre \emph{orienté} associé à $M$ est défini par le torseur $P_M/\mathbb{F}_2^{\underline{\omega}}$ sous $M_0 = V_J(\mathbb{F}_2)$, i.e. par

\page{36}\note{p.~18 de l'auteur.}
$F(\Delta)/\{1, \underline{a}_\Delta\}$ (cf.\ (36)), qui est donc l'ens. des quatre bifaces de $\Delta$. Donc c'est l'octaèdre orienté canoniquement associé à l'octaèdre $\Delta$, savoir
\[
  (45) \qquad
  \left\{
  \begin{aligned}
    \Delta^+ &= \operatorname{Or}(\Delta) \wedge^{\{1, \underline{a}_\Delta\}} \Delta \\
    \Delta &= \operatorname{Or}(\Delta) \wedge^{\{1, \underline{a}_\Delta\}} \Delta^+
  \end{aligned}
  \right.
\]
dont on a \add{bien}, comme il se doit
\[
  (45) \qquad
  \left\{
  \begin{array}{l}
    \operatorname{Or}(\Delta^+) \simeq \{\pm 1\} = \mathbb{1}_{\mathbb{F}_2} \\
    \underline{\omega}_{\Delta^+} \simeq \underline{\omega}_\Delta \simeq \underline{\omega}_J = \underline{\omega} \quad [\text{car } \underline{a}_\Delta \text{ opère trivialement sur } \underline{\omega}_\Delta] \\
    C_{\Delta^+} \simeq C_\Delta \wedge \mu \simeq \underline{\omega} \qquad \ldots
  \end{array}
  \right.
\]
\note{dans la première ligne de (45), un symbole noirci précède $\operatorname{Or}(\Delta)$, et l'indice $\{1, \underline{a}_\Delta\}$ est écrit à côté d'un premier indice biffé. Le numéro (45) est porté deux fois sur la page, ici et à la formule précédente ; transcrit tel quel (la suite de la page cite (47) et (48)).}

$[\![$ Il importe maintenant de bien comprendre la structure d'extension
\[
  (47) \qquad 1 \to \underset{\substack{\wr\\ \{\pm 1\}}}{\mathfrak{z}} \longrightarrow \operatorname{GL}(M) \longrightarrow \underset{\substack{\wr\\ \operatorname{Aut}_{\mathrm{oct}}(\Delta_M)'\\ \simeq \operatorname{Aut}^+_{\mathrm{oct}}(\Delta_M) \times \{1, \underline{a}_\Delta\}}}{\operatorname{GL}(M)'} \longrightarrow 1
\]
et en particulier
\[
  (48) \qquad 1 \to \underset{\substack{\|\\ \{\pm 1\}}}{\mathfrak{z}_0} \longrightarrow \operatorname{GL}(2, \mathbb{Z}/4\mathbb{Z}) \longrightarrow \underset{\substack{\wr\\ \operatorname{Aut}_{\mathrm{oct}}(\Delta_{(\mathbb{Z}/4\mathbb{Z})^2})\\ \text{« octaèdre standard »}}}{\operatorname{GL}(2, \mathbb{Z}/4\mathbb{Z})'} \longrightarrow 1
\]
puisque c'est par la cohomologie de cette extension que l'on comprend quelle est la donnée supplémentaire qu'il faut \uncertain{pour décrire}, d'un octaèdre non orienté $\Delta$, les $(\mathbb{Z}/4\mathbb{Z})$-modules libres de rang $2$ $M$ dont $\Delta$ proviendrait --- question qui revient à celle de « remonter » un torseur sous le groupe $\operatorname{GL}(2, \mathbb{Z}/4\mathbb{Z})'$ en un torseur sous le groupe $\operatorname{GL}(2, \mathbb{Z}/4\mathbb{Z})$ lui-même.
\note{le crochet ouvrant double, en tête de paragraphe, n'est pas refermé sur la page. Le prime de $\operatorname{Aut}_{\mathrm{oct}}(\Delta_M)'$ dans (47) est transcrit tel quel. La première partie de la dernière phrase (« puisque c'est par la cohomologie de cette extension que l'on comprend ») est d'une lecture incertaine.}
\page{37}
J'en reviens au diagramme $(23)$

\begin{tikzcd}[column sep=small, row sep=small]
    & \mathbb{F}_2^{\underline{\omega}} \arrow[dl] \arrow[dr, Rightarrow] & \\
    \operatorname{GL}_\Lambda(M) \arrow[d] \arrow[dr] & & \mathbb{F}_2^{\underline{\omega}} \simeq V(S_Q) \arrow[dl] \\
    \mathfrak{S}_I \arrow[dr] & \mathfrak{D}_Q \arrow[d, "\chi_C"] & \\
    & \{\pm 1\} &
  \end{tikzcd}

qui exprime $\operatorname{GL}_\Lambda(M)$ comme produit fibré, la donnée de $M$ équivalant à celle des triples $(I, Q, K)$, où $I$ est un ens. à quatre éléments, $Q$ un cube \add{combinatoire}, et $K$ un isom.
\[
  K : \underset{\text{codiagonales}}{C_Q} \xrightarrow{\;\sim\;} \underline{\omega}_I
\]
\note{diagramme redessiné : le carré $\operatorname{GL}_\Lambda(M)$, $\mathfrak{S}_I$, $\mathfrak{D}_Q$, $\{\pm 1\}$ porte au centre la mention « cart » (cartésien) ; la flèche de $\mathfrak{S}_I$ vers $\{\pm 1\}$ n'a pas d'étiquette. La double flèche du haut est tracée comme telle. Sous $\{\pm 1\}$, un « $\mathfrak{z}$ » (?) est griffonné. La flèche de $\mathbb{F}_2^{\underline{\omega}}$ au sommet vient d'un trait qui part du bord supérieur de la page.}

En termes de $Q$, on a
\[
  C_\Delta \simeq F(\Delta)/M_0 \simeq P_{M, \mathbb{F}_2^{\underline{\omega}}}/\mathbb{F}_2 = S_Q/\mathbb{F}_2 \simeq D_Q
\]
ens. des diagonales de $Q$ --- donc la donnée de l'octaèdre $\Delta$ (en termes de $M \leftrightarrow (I, Q, K)$) équivaut à celle de $I$, et des deux ens. à deux éléments $C$ $(\simeq \underline{\omega}_I)$ et $D$ correspondant aux diagonales et codiagonales du cube $Q$, alors que la donnée équivaut à celle d'un torseur sur $\mathbb{D}_4$ \uncertain{$\mathfrak{Z}(\mathbb{D}_4)$} $\simeq \mathbb{F}_2 \times \mathbb{F}_2$, qu'il faut donc relever en un torseur sous $\mathbb{D}_4$. Ou encore, décrire la donnée d'un cube combinatoire, dont on a déjà les deux ensembles \add{$C, D$} des codiagonales et diagonales\ldots\ $]\!]$
\note{le double crochet fermant clôt celui ouvert page~36. Devant $\mathbb{D}_4$, un symbole biffé ; la lecture de « $\mathfrak{Z}(\mathbb{D}_4)$ » (sans doute le quotient de $\mathbb{D}_4$ par son centre, ou son abélianisé) n'est pas sûre, et la phrase « alors que la donnée équivaut à celle d'un torseur sur\ldots » ne se construit pas tout à fait ; transcrite telle quelle.}
\subsection{Récapitulation partielle}

\page{38}
À un module $M$ libre de rang $2$ sur $\mathbb{Z}/4\mathbb{Z}$, on associe
\[
  (49) \qquad
  \left|
  \begin{array}{l}
    \text{a) L'octaèdre (non orienté) } \Delta_M = M^{\#}/\pm 1 \ (\text{ens. des sommets}), \\
    \quad \text{l'ens. des trois diagonales s'identifiant à } (M/2M)^* \simeq M_0^{\#} = J \\
    \text{b) L'ensemble à quatre éléments } I, \text{ qu'on peut p.\ ex.\ déduire de } \Delta \\
    \quad \text{comme } I_M \simeq \underset{\text{faces de } \Delta}{\underbrace{F(\Delta)}}/\text{antipodisme} \\
    \text{c) Le cube } Q_M, \ S_M = P_M/M_0
  \end{array}
  \right.
\]
où $P_M$ est le $\operatorname{End}(M_0)$-torseur formé des sections de $M^{\#}$ sur $M_0^{\#}$ satisfaisant $\xi_i + \xi_j + \xi_k = 0$, et où on identifie $\operatorname{End}(M_0)$ à $M_0 \oplus \mathbb{F}_2^{\underline{\omega}}$, de sorte que $P_M/M_0$ est un $\mathbb{F}_2^{\underline{\omega}}$-torseur, donc est l'ens. des sommets d'un cube \add{$Q = Q_M$} dont l'ens. des codiagonales est $\underline{\omega}$
\marginal{On a un iso. can. $J \simeq \mathfrak{P}_{2,2}(I_M)$, $V_J(\mathbb{F}_2) \simeq M_0$ $(\simeq \mathbb{F}_2^J/\mathbb{F}_2)$, $I_M$ torseur sous, $I_M \simeq P_M/\mathbb{F}_2^{\underline{\omega}}$}
\marginal{où $\underline{\omega} \simeq \underline{\omega}_I \simeq \underline{\omega}_J$}
\marginal{\emph{NB} $F(\Delta) \simeq I \times C_\Delta$ (« couleurs » des faces)}
\note{les trois points a), b), c) sont réunis par une accolade sous le numéro (49). Sous b), entre « faces de $\Delta$ » et c), une ligne interlinéaire en partie biffée (« ens.\ des sections de $\Delta$ \ldots\ ens.\ de sommets ») n'est pas reprise. En c), le symbole $S_M$ est écrit sur un signe noirci. Les deux premières notes marginales sont dans la marge gauche ; la troisième, encadrée d'un trait, à droite de b). Dans la première, « $I_M$ torseur sous » n'est pas suivi du groupe sur la page.}

Les données a) et b) ne dépendent que de la droite projective $X_M = \mathbb{P}(M)$ sur $\mathbb{Z}/4\mathbb{Z}$ ($\Delta$ est l'ens. des pts de celle-ci, $J$ l'ens. des pts de $(X_M)_0$, droite projective sur $\mathbb{F}_2$), mais non pas $Q_M$. On a des bijections canoniques
\[
  (50) \qquad
  \left\{
  \begin{array}{llll}
    \underset{\text{codiag.\ de } Q}{C_Q} \simeq \dot{\underline{\omega}}_\Delta & \overset{\mathrm{déf}}{\simeq} \underline{\omega}_I \simeq \underline{\omega}_J & & \text{soit } \underline{\omega} \\
    \underset{\text{diag.\ de } Q}{D_Q} \simeq \underset{\text{« couleurs » des faces}}{C_\Delta} & \simeq \underline{\omega} \wedge \mu & & \text{où } \mu \text{ défini ci-dessous} \\
    \underset{\text{or.\ de } Q}{\underline{\omega}_Q} \simeq \operatorname{Or}(\Delta) & \simeq \det(M)^* & & \text{soit } \mu
  \end{array}
  \right.
\]
\note{le tableau (50) est redessiné. À la première ligne, sous $\dot{\underline{\omega}}_\Delta$, un passage de deux lignes est biffé de hachures et n'est pas lu. Le point sur $\underline{\omega}_\Delta$ est de l'auteur. Au début de la phrase « Les données\ldots », « fibré » (?) est biffé devant « droite projective ».}
\page{39}
\note{feuillet écrit dans l'autre sens (en largeur), d'une encre plus pâle et d'une écriture plus grosse que les pages qui précèdent ; les ensembles y sont soulignés ($\underline{I}$, $\underline{J}$, $\underline{K}$), comme des objets au-dessus d'une base $S$. Le lien avec la récapitulation de la page~38 n'est pas explicite.}

$\underline{I}$ \ill{}

$(X, \underline{I})$ \struck{\ill{}} \ill{} \ill{} \ill{} \ill{} ?
\[
  \begin{array}{c}
    \big\updownarrow \\
    \bigl(\underline{I},\ \text{d'où } \underline{J},\ t \in \Gamma(\Sigma_J^*/S)\bigr) \longleftrightarrow \bigl(\underline{J}\ (\text{d'où } \underline{V}_J \text{ et } \Sigma_J),\ \text{section de } \Sigma_J^* \text{ sur } S,\ \text{torseur } I \text{ sous } \underline{V}_J\bigr) \\
    \big\updownarrow \\
    \bigl(\underline{J},\ t \in \Gamma(\Sigma_J^*/S),\ X \text{ rev.\ de } \Sigma_J^*,\ \text{étale sur } \Sigma_J^*,\ \ldots\bigr) \\
    \big\updownarrow \\
    \bigl(\underline{J},\ t \in \Gamma(\Sigma_J^*/S),\ \underline{K} \text{ torseur sous } \underline{J}\bigr)
  \end{array}
\]
\emph{NB} $t$ et $\underline{I}$ indépendants\ldots

\emph{NB} $X$ et $t$ sont indépendants l'un de l'autre.

\emph{NB} $\underline{K}$ et $t$ sont indépendants l'un de l'autre.

\emph{NB} $X \mid \underline{J} \simeq \Delta$ est la configuration octaédrale\ldots
\note{entre la première flèche et la ligne $(\underline{I}, \ldots)$, une ligne entière est biffée : \struck{$(\underline{J}(\Sigma_J), t \in \Gamma(\Sigma_J^*/S), \underline{I}$ \ill{} torseur sous}. La troisième ligne se poursuit, sous « $X$ rev.\ de $\Sigma_J^*$, étale sur $\Sigma_J^*$ », par « galoisien sur \ldots\ fibré \ldots\ rev.\ \emph{abélien} » et un passage biffé, d'une lecture très incertaine. Dans la marge droite, écrite verticalement, une suite d'inclusions ($\ldots \subset \ldots$, $\mathfrak{S}(M_0)$, $\operatorname{End}(M_0)$, « $1 \in$ ») reliée par une flèche au texte, et une petite figure en croix portant les lettres $x$, $X$ ; elle n'est pas lue. Le symbole biffé devant $\underline{K}$ n'est pas lu.}

Quand $\underline{I}$ est muni d'une section, de sorte que $\underline{I} \simeq \underline{J} \amalg S$, alors on a un iso. can. $X \simeq \Sigma_J$, d'où $f : \Sigma_J \to \Sigma_J$
\note{la page s'arrête sur cette formule ; l'indice de $f$ est noirci.}

\end{document}
